\subsection{The setting}
\label{subsec:setting}

Let $\widehat X$ be a smooth projective Calabi--Yau threefold and $F$ an open face of its nef
cone, that is, the relative interior of a face. We assume:
\begin{enumerate}
\item[(S1)] for $J\in F$ one has $J^3>0$, the irreducible curves $C$ with $J\cdot C=0$ are exactly $k$ pairwise
 disjoint $(-1,-1)$-curves $C_1,\dots,C_k$, and $F$ is open in the annihilator
 $\operatorname{ann}(C_1,\dots,C_k)\subset H^2(\widehat X;\mathbb R)$ of their classes;
\item[(S2)] the classes of the $C_i$ satisfy exactly one relation $\sum_ia_i[C_i]=0$, with $a$
 primitive and every $a_i\neq0$, and the lattice $\Lambda\subset\ZZ^k$ of charge vectors
 $(D\cdot C_1,\dots,D\cdot C_k)$, $D\in H^2(\widehat X;\ZZ)$, is saturated, so that
 $\Lambda=a^\perp\cap\ZZ^k$ has rank $k-1$;
\item[(S3)] the gauge kinetic matrix of the vector multiplets, the Hodge metric
 $\int\omega\wedge*_J\omega'$ on harmonic two-forms, is positive definite on $F$.
\end{enumerate}
We call a face satisfying (S1)--(S3) a \emph{face with one relation}. On $F$ the Hodge metric is a limit of Hodge metrics on the K\"ahler cone, hence positive
semidefinite, so (S3) is equivalent to its nondegeneracy there. Saturation in (S2) only
normalises the gauge torus: in general the torus acts on the $k$ hypermultiplets below through
its image, whose cocharacter lattice is $a^\perp\cap\ZZ^k$, and a finite subgroup acting
trivially on them does not change the germs studied here.

Under (S1) a rational class in $F$ is nef and big and defines a contraction
$\pi_F:\widehat X\to X_F$ with $k$ ordinary double points, the images of the $C_i$. Since all
$a_i$ are nonzero, $X_F$ has a first-order smoothing by Friedman's
criterion~\cite{Friedman:1986}, its deformations are unobstructed~\cite{Gross:1997}, and its smoothings $X_t$ satisfy
$h^{1,1}(X_t)=h^{1,1}(\widehat X)-(k-1)$ and $h^{2,1}(X_t)=h^{2,1}(\widehat X)+1$. In M-theory on
$\widehat X$ with K\"ahler class $J\in F$, M2-branes on the $C_i$ give $k$ massless
hypermultiplets, one for each curve: an isolated $(-1,-1)$-curve has genus-zero
Gopakumar--Vafa invariant $1$ and no multiple-cover or higher-genus
contributions~\cite{FaberPandharipande:2000,BryanKatzLeung:2001}. They carry charges in
$\Lambda$ under the torus $T^{k-1}$ that is the image of $H^2(\widehat X;\mathbb R)/H^2(\widehat
X;\ZZ)$ in $U(1)^k$. The $k$ flop walls through $F$ make the cubic prepotential of $\widehat X$
only twice continuously differentiable along $F$. Writing $Z_i=J\cdot C_i$ for the central charges of the light states and
keeping them as fields, the prepotential of the Wilsonian theory in the chamber where all
$Z_i>0$ is
\begin{equation}
 \mathcal F_W=\mathcal F_{\widehat X}+\tfrac1{12}\textstyle\sum_iZ_i^3,
 \label{eq:wilsonian-F}
\end{equation}
the Chern--Simons levels shifting by $\frac12q\otimes q\otimes q$ for each hypermultiplet of
charge $q$~\cite{Witten:1996qb,Intriligator:1997}. It is a polynomial with the same Hessian as
$\mathcal F_{\widehat X}$ at $Z=0$.

For $X_9$, Propositions~\ref{prop:face} and~\ref{prop:face-smoothing} of
Section~\ref{sec:face} show that $\widehat X=\wideX_9$ and the face~\eqref{eq:face-def} satisfy
(S1)--(S3) with $k=4$ and $a=(1,1,-1,1)$; see the proof of Theorem~\ref{thm:four-node}.

\subsection{Gauged supergravity and its Minkowski vacua}
\label{subsec:sugra}

M-theory on $\widehat X$ gives five-dimensional $N=1$ supergravity with
$n_V=h^{1,1}(\widehat X)-1$ vector multiplets and $n_H=h^{2,1}(\widehat X)+1$ neutral
hypermultiplets; for $\wideX_9$ these are $5$ and $123$. The vector-multiplet scalars are the
very special real coordinates $h^I$, $I=1,\dots,n_V+1$, on the hypersurface
$\frac16\kappa_{IKM}h^Ih^Kh^M=1$~\cite{Gunaydin:1984}; the overall volume of $\widehat X$ sits
in the universal hypermultiplet. The hypermultiplet scalars take values in a quaternionic
K\"ahler manifold $M$ of negative scalar curvature. Supergravity fixes only the reduced scalar
curvature of $M$, in units of the five-dimensional Planck mass; the hyperk\"ahler part of the
curvature is not fixed. We measure distances on $M$ in Planck units. The transverse
correction of order $r^2$ found below is given by the curvature of $M$ at the root
(Proposition~\ref{prop:universal}), which may be large where the neutral moduli approach a degeneration.
At a point where charged states become massless the effective theory includes them as further
hypermultiplets, and the vectors gauge isometries of $M$ with Killing vectors $k_I$ and
quaternionic moment maps $P_I$; up to normalisation these are the moment maps $\mu_\xi$ of
Section~\ref{sec:quotients} for $\xi$ the generator dual to the $I$-th basis divisor.

For constant scalars and vanishing field strengths, the fermion shifts of the gravitino,
gaugini and hyperini give the conditions~\cite[eq.~(2.9)]{W26quantum}; compare the ground-state
conditions of~\cite[eqs.~(5.17)--(5.18)]{Ceresole:2000}:
\begin{equation}
 h^IP_I=0,\qquad h^I_xP_I=0,\qquad h^Ik_I=0 ,
 \label{eq:shifts}
\end{equation}
where $h^I_x$ are the tangent vectors to the very special real hypersurface. The gaugino shift
is proportional to $g^{xy}h^I_yP_I$, with $g_{xy}$ the vector-multiplet metric; where $g_{xy}$ is
nondegenerate its vanishing is equivalent to the second condition. The vectors $h^I$ and $h^I_x$
always span $\mathbb R^{n_V+1}$, so where the vector-multiplet metric is nondegenerate the
conditions become
\begin{equation}
 P_I=0\quad\text{for all }I,\qquad h^Ik_I=0
 \label{eq:vacua}
\end{equation}
\cite[eq.~(2.10)]{W26quantum}; the four-dimensional analogue is~\cite[Sect.~2.4.2]{Hristov:2009}.
The potential vanishes on these solutions. The action of M-theory also contains
higher-derivative terms, among them the supersymmetric completion of
$A\wedge\operatorname{tr}R\wedge R$~\cite{Hanaki:2006}. That they do not change the set of such
vacua is part of Hypothesis~\ref{hyp:regular} below. It is supported by the observation that on
a background with constant scalars and vanishing field strengths the auxiliary fields of the
Weyl multiplet are set to zero by the supersymmetry variations and by Lorentz invariance, so
that the auxiliary equation of the vector multiplets still reduces to the vanishing of the
moment maps.

\subsection{The local vacuum space}
\label{subsec:object}

We write $N$ for the quaternionic K\"ahler manifold of neutral hypermultiplets of $\widehat X$,
of quaternionic dimension $h^{2,1}(\widehat X)+1$, and call $b\in N$ \emph{smooth} if $N$ is a
smooth manifold near $b$, that is, if $b$ lies over a point of the complex-structure moduli space
of $\widehat X$ away from the loci where $\widehat X$ acquires further singularities.

\begin{definition}\label{def:vacuum-germ}
Let $\phi\in F$ be normalised by $J^3=6$ and $b\in N$ a smooth point. The \emph{root point}
$x=(\phi,b,0)$ is the vacuum with vector moduli $\phi$, neutral hypermultiplets $b$ and
vanishing charged hypermultiplets. The \emph{local vacuum space} $\mathcal V_x$ is the germ at
$x$ of the set of constant-scalar, zero-field-strength solutions of~\eqref{eq:vacua} in the
theory of Hypothesis~\ref{hyp:regular} below, modulo gauge equivalence, regarded as a
stratified topological space.
\end{definition}

The \emph{Coulomb branch} through $x$ is the part of $\mathcal V_x$ on which the charged
hypermultiplets vanish; we write $\mathcal K_\phi$ for its vector-multiplet germ at $\phi$. It is
the germ at $\phi$ of the hypersurface of normalised vector-multiplet scalars, divided by the
walls $Z_i=0$ into chambers subject to $\sum_ia_iZ_i=0$; among these chambers are the K\"ahler
cone of $\widehat X$ and those of its flops along the $C_i$. The \emph{Higgs branch} is the closure of
the part of $\mathcal V_x$ on which the gauge group acting on the charged hypermultiplets is
broken to a finite group, and a \emph{mixed branch} would be a component on which some, but not
all, of the charged hypermultiplets are nonzero. We write $F_\phi$ for the germ at $\phi$ of
$F\cap\{J^3=6\}$. On the Higgs branch away from its root the compactification is M-theory on a
smoothing $X_t$, $t\neq0$: we call this part the \emph{regular Higgs branch}
$\mathcal H^{\mathrm{reg}}$. At each of its points the vector moduli are normalised K\"ahler
classes of $X_t$, the hypermultiplet moduli form a smooth quaternionic K\"ahler manifold
$\mathcal M_H^{\mathrm{reg}}$ of quaternionic dimension $h^{2,1}(X_t)+1$, and the scalar manifold
is locally the product of the two. At a root where no Lagrangian description of the light
states is available, the Higgs branch is defined through $\mathcal H^{\mathrm{reg}}$ alone: it is
the metric completion of $\mathcal H^{\mathrm{reg}}$ near the limit point, for the metric of
$\mathcal M_H^{\mathrm{reg}}$ on the hypermultiplet factor and the very special real metric on
the vector factor.

\subsection{The regularity hypothesis}
\label{subsec:hypotheses}

\begin{hypothesis}[regularity]\label{hyp:regular}
Near every point of $F$ and every smooth point of $N$, the dynamics below the lightest massive
state is five-dimensional $N=1$ gauged supergravity with: the vector multiplets of $\widehat X$
with the Wilsonian prepotential~\eqref{eq:wilsonian-F}; a smooth quaternionic K\"ahler manifold
$M$ of quaternionic dimension $h^{2,1}(\widehat X)+1+k$, containing $N$ and the $k$ charged
hypermultiplets; and the gauging of the torus $T^{k-1}$ of isometries of $M$ with charge
lattice $\Lambda$, acting on a neighbourhood of $N$. No neutral hypermultiplet is charged.
Moreover:
\begin{enumerate}
\item the Coulomb branch through each root point, M-theory on $\widehat X$ with K\"ahler class
 near $F$ and vanishing charged hypermultiplets, consists of supersymmetric Minkowski vacua of
 this theory, and higher-derivative terms do not change the set of such vacua;
\item the K\"ahler cone of $X_t$ is the same for all smoothings $X_t$ parametrised by a
 neighbourhood $W$ of each root point $x$ in the hypermultiplet directions, and at every point of
 $W\cap\mathcal M_H^{\mathrm{reg}}$ the vector moduli of the regular Higgs branch form the whole
 normalised K\"ahler cone $\mathcal K(X_t)$;
\item near each root point, the Higgs vacua of this theory off the root are M-theory on the
 smoothings $X_t$, and $\mathcal M_H^{\mathrm{reg}}$ is bi-Lipschitz equivalent, near the root,
 to $(\mu^{-1}(0)\smallsetminus N)/T^{k-1}$ with the quotient metric of~\cite{Galicki:1988}.
\end{enumerate}
\end{hypothesis}

For $X_9$ the hypothesis concerns $M$ of quaternionic dimension $127$, the four charged
hypermultiplets of Proposition~\ref{prop:face} and the torus $T^3$ with the
charges~\eqref{eq:app-Q}.

The hypothesis has the form of the standard treatment of conifold singularities. In four
dimensions the singularity of the vector-multiplet moduli space at a conifold point is an
effect of integrating out the light states: the Wilsonian action that keeps them is
smooth~\cite{Strominger:1995Conifold,Greene:1995hu}. On the hypermultiplet side, for type~IIA
string theory near $k$ homologous vanishing three-cycles, the metric on the neutral fields alone,
corrected by Euclidean D2-branes on the vanishing cycles, has a $\CC^2/\ZZ_k$
singularity~\cite{OoguriVafa:1996,SeibergShenker:1996}, the same as the hyperk\"ahler quotient of
the Wilsonian description with $k$ light hypermultiplets. The hypothesis asserts the existence
of the corresponding smooth Wilsonian hypermultiplet manifold in five dimensions at finite Planck
mass, as in the supergravity models of~\cite{Mohaupt:2004de}; its metric is not
known~\cite{Mohaupt:2004de}, and nothing below depends on it. The geometric input is proved for $X_9$ in Proposition~\ref{prop:face}: exactly four hypermultiplets are massless, no
divisor collapses, and all gauge couplings are finite. The torus acts on a neighbourhood of $N$
because its weights are integral.

The existence of $M$ and clause (1) are the standard assumptions of the conifold analyses cited
above, carried over to finite Planck mass; they are all that Theorem~\ref{thm:general} uses.
Clauses (2) and (3) are technical: they are used only in Lemmas~\ref{lem:product}
and~\ref{lem:completion} and Theorem~\ref{thm:wall}, to pass from the four-node roots of $X_9$ to
the other roots through metric completions, which bi-Lipschitz equivalences preserve.

In detail, clause (1) says that the Wilsonian theory reproduces the
supersymmetric compactifications of M-theory on $\widehat X$, including those with the K\"ahler
class near $F$, and that the four-derivative terms do not alter this; it is what forces the
moment map to vanish at the root. Clause (2) excludes a dependence of the chamber structure on
the complex structure of $X_t$ near the root. Such a dependence does occur in general: the
K\"ahler cone of a Calabi--Yau threefold can jump under deformation of its complex structure,
and by Wilson's theorem this happens only on members containing elliptic quasi-ruled
surfaces~\cite{Wilson:1992,Wilson:1993}. So the first half of clause (2) holds if no $X_t$
parametrised by $W$ contains such a surface; the second half, that the vector moduli of the
regular Higgs branch fill the whole normalised K\"ahler cone over $W$, is an independent
assumption. Clause (3) identifies distances measured on the smooth side with those of the
Wilsonian quotient up to bounded distortion.

\begin{lemma}\label{lem:product}
Over the neighbourhood $W$ of clause (2) of Hypothesis~\ref{hyp:regular}, the metric of
$\mathcal M_H^{\mathrm{reg}}$ on the regular Higgs branch does not depend on the vector-multiplet
moduli.
\end{lemma}

\begin{proof}
At every point of $\mathcal H^{\mathrm{reg}}$ the compactification is M-theory on a smooth
threefold $X_t$ with K\"ahler class in the interior of its K\"ahler cone, described by
two-derivative supergravity with neutral hypermultiplets. There supersymmetry forbids neutral
couplings between the two kinds of multiplets at the two-derivative level, and the scalar
manifold is locally the product of the very special real and the quaternionic K\"ahler
factors~\cite[Sect.~2.4]{Becker:1995kb}, the form assumed in~\cite{Ceresole:2000} and the same
factorisation that identifies the hypermultiplet moduli spaces of M-theory and of type~IIA string
theory on one threefold~\cite[Sect.~2]{Strominger:1997}. So the hypermultiplet metric is locally
constant in the vector moduli. By clause (2) of Hypothesis~\ref{hyp:regular}, over $W$ the
regular Higgs branch is, as a set, the product of $W\cap\mathcal M_H^{\mathrm{reg}}$ with the
whole normalised K\"ahler cone of $X_t$; the latter is identified by radial projection with a convex set, hence
connected, and the metric is therefore independent of the vector moduli. The moduli space of
vacua does not factorise at transition points themselves~\cite[Sect.~2.1]{Aspinwall:2000}, but
$\mathcal H^{\mathrm{reg}}$ contains none.
\end{proof}
